% Part I, chapter 4: Differences from the original package.
\chapter{Differences from the original package}
\label{chap:behavior}

\wfdbx{} reproduces the WFDB Software Package \wfdbversion{} exactly, with
the deliberate exceptions described in this chapter. It has no built-in
signal parameters; it finds files only where they are named, and opens URLs
only when network access is on; it sorts out-of-order annotation files
itself, without losing annotations; it never writes a signal file over a
signal file it is reading; its tools report write errors in their exit
status; its \code{getinfo} finds the record line of a header as the header
parser does; and it gives a defined result for input on which the C package
crashes. It also adds features that the C package does not have.

These exceptions are the same in all seven implementations, and are tested
like everything else: the reference results of the conformance suite come
from the C package with the patch \file{tests/reference/divergences.patch}
applied, which is the precise definition of the differences. The original
behavior remains available in every case, on request, in \emph{legacy mode}
(\cref{sec:legacy}). The file \file{docs/DIVERGENCES.md} gives the reasons
for each difference in more detail, and \cref{sec:corrections-1-2-0,sec:corrections-1-2-2} list
the corrections of versions 1.2.0 and 1.2.2.

\section{No built-in sampling frequency, gain, ADC resolution or units}
\label{sec:no-defaults}

The C package falls back on built-in values for four parameters of a
signal whenever a record or the user does not give them
(\cref{tab:c-defaults}). These values describe the equipment with which the
first PhysioNet databases were recorded, and do not describe modern
recordings. A silently assumed value corrupts every time, interval,
resampled signal and physical value computed from it, and nothing in the
output says that it was assumed. \wfdbx{} therefore never assumes any of
these four parameters: it uses the values given in the record, or values
given explicitly by the user, and reports an error when an operation needs a
value that is not available. Where an output must still show missing units,
it shows \code{nd} (``not defined'', a notation that the C package already
used for empty units) rather than \code{mV}.

\begin{table}[htbp]
\caption{The values that the C package assumes.}
\label{tab:c-defaults}
\begin{tabularx}{\linewidth}{@{}P{3cm}L@{}}
\toprule
{\tabhead Parameter} & {\tabhead Built-in value in C} \\
\midrule
sampling frequency & 250\,Hz (\code{WFDB\_DEFFREQ}); 500\,Hz in
  \cmd{parsescp} \\
gain & 200\,adu/mV (\code{WFDB\_DEFGAIN}); 400\,adu/mV for AHA files; 12.5 or
  6400\,adu/mV for WAV files \\
ADC resolution & 12 bits (\code{WFDB\_DEFRES}), or the storage width of the
  format \\
units & millivolts \\
\bottomrule
\end{tabularx}
\end{table}

\paragraph{Library.}
\begin{itemize}
\item A header that does not give a sampling frequency leaves it
  unspecified: \code{sampfreq} returns 0 (C: 250), unless a frequency was set
  with \code{setsampfreq} before the record was opened (as in C).
\item A header that does not give an ADC resolution leaves it unspecified
  (0); C inferred it from the format or assumed 12 bits. Missing units remain
  unspecified (\code{NULL}), as in C, but are never taken to be millivolts.
\item The time conversions that need the frequency fail when it is
  unspecified: \code{strtim} reports \code{strtim: sampling frequency is not
  specified} and returns 0 for the time strings other than sample numbers
  (\code{s\ldots}), \code{e}, \code{i} and \code{o} (and counter values when
  there is no counter frequency); \code{timstr} and \code{mstimstr} report
  the same error and return the time as a sample number (\code{s} followed
  by the number). C used 1 sample per second. The new function
  \code{wfdbx\_time\_needs\_freq(string)} tells whether \code{strtim} needs
  the frequency for a string (\cref{sec:times}).
\item \code{setifreq} fails if the input frequency is unspecified
  (\code{setifreq: sampling frequency is not specified}).
\item \code{aduphys}, \code{physadu}, \code{adumuv} and \code{muvadu} fail
  for an uncalibrated signal (gain 0) or an invalid signal number: they
  report the error and return NaN (\code{aduphys}) or
  \code{WFDB\_INVALID\_SAMPLE} (the others). C assumed a gain of 200, and a
  baseline of 0 for invalid signal numbers.
\item The constants \code{WFDB\_DEFFREQ}, \code{WFDB\_DEFGAIN} and
  \code{WFDB\_DEFRES} remain defined, for compatibility, but the library does
  not use them (except in legacy mode).
\end{itemize}

\paragraph{Tools.} A tool that writes these parameters into a record
requires them; a tool that needs a parameter missing from its input stops
with an error only when the requested operation needs it:
\begin{itemize}
\item \cmd{wrsamp} requires \opt{-F}, \opt{-G} and the new option
  \opt{-R BITS} when it writes a record (\opt{-o}), and the new option
  \opt{-U UNITS} for calibrated signals unless the input gives the units;
  \opt{-G 0} declares the signals uncalibrated. \opt{-G}, \opt{-R} and
  \opt{-U} take one value per signal (the last applies to the rest).
\item \cmd{wav2mit} has the new options \opt{-G GAIN} (required) and
  \opt{-U UNITS} (required unless the gain is 0); \cmd{ad2m} and
  \cmd{ahaecg2mit} the new options \opt{-F FREQ}, \opt{-G GAIN},
  \opt{-R BITS} and \opt{-U UNITS}, all required (\opt{-U} unless the gain
  is 0).
\item \cmd{a2m}, \cmd{rr2ann}, \cmd{rdann}, \cmd{ann2rr}, \cmd{sortann} and
  \cmd{m2a} take \opt{-F FREQ} (new except in \cmd{rr2ann}), used when the
  record has no header or its header has no frequency, and required then
  when times are converted.
\item \cmd{rdsamp} needs the frequency for time arguments, time columns and
  XML output, and the gain of each selected signal for physical units
  (\opt{-p}); \cmd{sigamp}, \cmd{snip}, \cmd{wfdb2mat}, \cmd{mit2wav},
  \cmd{wfdbtime}, \cmd{xform} and \cmd{mit2edf} need the frequency (and
  \cmd{mit2edf} the gain and ADC resolution of every signal), and so does
  \cmd{time2sec} with \opt{-r} (since 1.2.2; C converted the time at
  250\,Hz, and \wfdbx{} until 1.2.1 printed a meaningless number,
  \code{-9223372036854775808}, with the exit status 0).
\item \cmd{wfdbdesc} prints \code{not specified} for an unspecified
  frequency and ADC resolution, \code{uncalibrated} for a gain of 0 (C:
  \code{uncalibrated; assume 200}), and \code{adu/nd} for unspecified units;
  \cmd{sampfreq} prints 0 for a record whose header gives no frequency (C:
  250); \cmd{parsescp} writes the frequency given in the SCP file (C always
  wrote 500\,Hz).
\end{itemize}
The messages are \code{TOOL: sampling frequency of record RECORD is not
specified} (followed by \code{ (use -F)} for the tools that have that
option), \code{TOOL: signal N of record RECORD is uncalibrated},
\code{TOOL: ADC resolution of signal N of record RECORD is not specified},
and \code{TOOL: the \ldots{} must be specified (-X)} for a missing required
option. \Cref{part:reference} gives the options of each tool.

\paragraph{High-level API.} \code{write\_record} requires the sampling
frequency (greater than 0), the gains, the ADC resolutions, and the units of
the calibrated signals; a gain of 0 declares a signal uncalibrated, which is
allowed only for samples in ADC units. \code{read\_record} reports an
unspecified frequency as 0, and fails for time strings that need it; with
physical units, it fails if a selected signal is uncalibrated.
\code{read\_annotations} takes an optional sampling frequency, used, like
the \opt{-F} option of the tools, when the record has no header or its
header gives none.

\paragraph{Reproducing the results of the original package.} The results of
the original tools are obtained by giving the values they assumed,
explicitly. For example, \cmd{wrsamp} writes a record as the original did
with \opt{-F 250 -G 200 -U mV} and \opt{-R} with the storage width of the
output format (\opt{-R 12} for format 212, \opt{-R 16} for formats 16, 61,
160 and 516, \opt{-R 8} for 8, 80 and 508, \opt{-R 10} for 310 and 311,
\opt{-R 24} for 24 and 524, \opt{-R 32} for 32); \cmd{wav2mit} with
\opt{-G 12.5 -U mV} for 8-bit files and \opt{-G 6400 -U mV} for 16-bit
files; \cmd{ad2m} and \cmd{ahaecg2mit} with \opt{-F 250 -G 400 -R 12 -U mV};
the annotation tools with \opt{-F 250} for a record without a frequency; and
a library program calls \code{setsampfreq(250)} before it opens a record
whose header gives no frequency. One difference remains by construction: a
header written with \opt{-U mV} says \code{200/mV} where the C tools wrote
\code{200}.

\paragraph{What does not change.} Values defined by a data format, not by an
assumption about the signal (the record lengths and time offsets of the AHA
database formats, and \cmd{md2a}'s default length); values read from the
input (the sampling frequency and resolution of a WAV file, the parameters
of an EDF file); and analysis settings that are documented options, such as
the measurement windows of \cmd{sigamp}. \cmd{parsescp} writes a gain of
200\,adu/mV because it rescales the SCP samples to units of 5\,$\mu$V, so
that value is exact, not assumed.

\section{Files are found only where they are named}
\label{sec:named-files}

The C library found a file by searching the WFDB path, a list of places
tried in order until the file was found, by default the current directory,
\file{/usr/local/database} and, in builds with network support, PhysioNet.
That was convenient, but:
\begin{itemize}
\item a search opens the first copy it finds, so a stale or partial copy in
  a directory earlier on the path is read instead of the intended one, and
  a signal file missing next to its header is replaced by a file of the same
  name found anywhere else on the path;
\item whether a record is read from disk or downloaded, and from where,
  depends on the environment, not on the program, and a missing local file
  turns into a network request instead of an error;
\item for every record it read, the library also asked the last entry of the
  path, PhysioNet, for an optional \file{.info} file when there was no local
  one, with the full local path of an absolute record name in the request;
\item a header could make the library read any local file
  (\code{../../.ssh/id\_rsa}) or contact a third-party server, through the
  names of its signal files;
\item the default URL is no longer where PhysioNet serves its files.
\end{itemize}
\wfdbx{} reads and writes every file at the one location that its name
designates, as \cref{chap:files} describes: under the input or output
directory for a relative name, as given for an absolute path, a
\code{\textasciitilde} path or a URL, and relative to the header for the
files that a header names, which must stay in the header's directory. URLs
are opened only when network access is on, and only \code{http://} and
\code{https://} URLs can be read. A name containing \code{..} is
accepted. \code{setwfdb} and \code{getwfdb} set and return the input
directory, and the \code{WFDB} environment variable is used only in legacy
mode. Reading PhysioNet records must therefore be asked for explicitly
(\cref{sec:network}); in return, every program reads exactly the files it
names, the same on every machine.

The files that a tool opens itself, by the name given on its command line,
rather than through the library (the \opt{-i} file of \cmd{wrsamp}, the WAV
file of \cmd{mit2wav}, and the others of \cref{tab:tool-files}), follow the
same rules: the C tools opened them relative to the current directory, so
that, with an input or output directory set, a tool read its records from
one place and its own files from another. In \wfdbx{} such a file is found,
or created, where the library finds, or creates, a file of that name
(\cref{sec:tool-files}); a URL is not accepted for these files. In legacy
mode, the files that a tool reads itself are opened as named, and the
output directory, if one is set, is a plain prefix of the names of the
files that it creates itself.

\section{Out-of-order annotation files are sorted in the library}
\label{sec:sorting}

When an annotation file written by the library is not in canonical order
(\cref{sec:annotations}), the C library sorted it when the file was closed by
running the \cmd{sortann} program through the shell, unless the environment
variable \code{WFDBANNSORT} was 0. The result then depended on which
\cmd{sortann} was found on the \code{PATH}, and a record or annotator name
with a space or a shell character broke the command.

\wfdbx{} sorts the file itself, in the same process. The file it writes is
byte for byte the one that \cmd{sortann} writes, in the same order (by time,
then annotation number, then channel; a later annotation with the same
time, number and channel replaces the earlier one), with the same treatment
of \code{NOTQRS} annotations (they are dropped) and of the time resolution,
and the same encoding. The messages and status of the library are the same:
\code{Rearranging annotations for output annotator ANN ...done!}, or the
advice \code{Use the command: sortann -r REC -a ANN} when sorting is turned
off or fails; only the messages that the separate \cmd{sortann} printed
itself are gone. \code{WFDBANNSORT=0} turns sorting off, and the
\code{sort} option of the high-level \code{write\_annotations} decides for
its call alone (sort, or leave the file as written), whatever
\code{WFDBANNSORT} says.

\paragraph{Sorting does not lose annotations.} The original \cmd{sortann}
silently loses annotations: when an annotation at time 0, with annotation
number 0 and channel 0, arrives after later ones, it is dropped, and so are
the annotations later inserted directly behind it (input sample times 5, 0,
3 give a file that contains only 5), because the search for the insertion
point mistakes the empty head entry of its list for an annotation. \wfdbx{}
puts such an annotation at the beginning of the list, where it belongs (5,
0, 3 give 0, 3, 5). Nothing else about the sorting changes. This applies to
the \cmd{sortann} tool and to the sorting done by the library.

\begin{note}
In legacy mode, the library runs \cmd{sortann} as the original does (with
its own \cmd{sortann} tool, started with \opt{--legacy}, without a shell),
and both lose these annotations as the original does. Legacy mode should not
be used to sort annotation files unless reproducing the original results,
losses included, is the point: nothing in the output says that annotations
were dropped.
\end{note}

\section{An output signal file is never one of the input signal files}
\label{sec:output-input}

The files of records can be shared: the segments of a multi-segment record
are other records, so that \cmd{xform -i vlay -o fmt16} reads
\file{fmt16.dat} (a segment of \code{vlay}) and writes \file{fmt16.dat}
(the output record \code{fmt16}). The original opens the output file, which
truncates it, while still reading it: the input is destroyed, and what is
read depends on how much of it the C library had read ahead, so that the
result differs between platforms.

\wfdbx{} refuses to create an output signal file that is the same file as a
signal file of the record open for reading (the same file, not the same
name: links and different spellings of a path are recognized), with the
message \code{osigfopen: output file NAME is also an input file}
(\code{osigopen:} when the file is opened by \code{osigopen}); nothing is
written. The check is in the library, so every tool and every program that
uses the library functions is protected. The high-level \code{write\_record}
reads nothing: it closes any open record first, and simply replaces
existing files of the same name. The check covers signal files only: for
example, \cmd{snip -i 100s -n 100s} still writes the header of its output
record over that of its input record, in C and in \wfdbx{} alike; the output
record needs a name of its own.

\section{Write errors give a nonzero exit status}
\label{sec:write-errors}

The original tools report a file that cannot be written (a full disk, for
example) only through the library's messages, such as \code{osigclose:
write error in signal group 0} or \code{oannclose: write error on annotation
file atr}, and then exit with status 0; they say nothing at all when the
standard output cannot be written, or when a file that a tool writes itself
fails. The tools of \wfdbx{} follow the convention of the standard Unix
tools: when an output file or the standard output cannot be written, the
exit status is 1. The library's messages are unchanged; a standard output
that cannot be written is reported as \code{PROGRAM: write error on standard
output}, and a file that \cmd{snip}, \cmd{mit2edf} or \cmd{parsescp} writes
itself as \code{PROGRAM: write error on FILE}. A tool that exits with
another nonzero status for another reason keeps that status, and the tools
keep running after a write error, as the original ones do.

The library counts the output files that it could not write, flush or close;
\code{wfdbx\_write\_errors()} returns that count, and
\code{wfdbx\_note\_write\_error()} adds one, for a program that writes a file
itself.

\section{The record line of a header is found as the header parser finds it}
\label{sec:record-line-divergence}

The info strings of a header are its lines that begin with \code{\#} after
the record line (\cref{sec:info-strings}). The C package has two rules for
finding the record line. The header parser (\code{readheader}) takes the
first line whose first token (separated by spaces, tabs, carriage returns
and line feeds) does not begin with \code{\#}, so it skips empty and blank
lines as well as comment lines; but \code{getinfo} takes the first line that
does not begin with \code{\#}. For a header whose first line is empty or
blank, such as
\begin{lstlisting}[style=text]
(an empty line)
# Age: 40
rec 1 250 10
rec.dat 16 200 12 0 0 0 0 ECG
# Age: 63
\end{lstlisting}
the original \code{getinfo} takes the empty line for the record line, and
returns \code{Age: 40}, a comment before the record line, as an info string,
as well as \code{Age: 63}: \cmd{wfdbdesc} prints both as notes, and the
\code{comments} of the record hold both.

Since version 1.2.2, the \code{getinfo} of \wfdbx{} finds the record line as
\code{readheader} does, so that the info strings are the lines that the
header parser leaves after the record line: \code{Age: 63} only. A header
with no record line has no info strings. The info strings after the record
line are found as before: every line whose first character is \code{\#} (a
comment between two signal lines included). This applies to \code{getinfo}
itself, to the \code{comments} of \code{read\_header} and
\code{read\_record}, and to the notes of \cmd{wfdbdesc}. The metadata
functions (\code{read\_metadata} and the others, and the \code{metadata} of
\code{read\_header}) follow the same rule in every mode, legacy mode
included (\cref{chap:metadata}). In legacy mode, \code{getinfo} follows the
original rule. The conformance cases \code{edge-hdr-leading-*} check both (in
the legacy suite, against the unmodified package).

\section{Input on which the C package crashes}
\label{sec:crashes}

Some malformed or unusual input makes the C package crash, run out of
memory, loop forever, or give a result that depends on the machine (reads of
uninitialized memory). A crash cannot be reproduced, so \wfdbx{} gives each
such input one defined, safe result, the same in all seven implementations,
and legacy mode gives the same result. For example:
\begin{itemize}
\item a header with a negative or impossibly large number of signals
  (\code{n -1 250 10}), for which C tries to allocate memory for 4,294,967,295
  signals, gives \code{init: unexpected EOF in header file for record n}, as
  C does for a header with fewer signal lines than it declares;
\item an EDF file whose signals have 0 samples per data record, on which C
  crashes with a floating-point exception, reads as a record with no
  samples;
\item a FLAC signal file with 3 or more samples per frame whose FLAC block
  size is not a multiple of them, for which C writes past the end of a
  buffer, is read correctly;
\item \cmd{wfdb2mat} with a one-character record name, for which C
  overwrites the header of its input record, names its output record as for
  longer names (\code{mm} for \code{m});
\item \cmd{edf2mit} with a file shorter than the 256-byte fixed part of the
  EDF header prints \code{edf2mit: incomplete EDF header in FILE} and exits
  with status 2, creating no files;
\item \cmd{signame -s N}, where \code{N} is the number of signals, prints
  \code{[INVALID]}, as its manual page says;
\item \code{osigopen(record, NULL, 0)}, for which C reads before the start of
  an array, closes the output signals already open (unless the record name
  begins with \code{+}), reads the header (an error if it cannot be read),
  opens nothing and returns 0.
\end{itemize}
The complete list, with the C behavior and the \wfdbx{} choice for each, is
the table ``Undefined behaviors (not reproduced)'' of
\file{docs/KNOWN\_C\_BEHAVIORS.md}. The conformance suite checks most of
these choices, with reference results recorded from \wfdbx{} instead of C.

\section{Additions}
\label{sec:additions}

These are new features, not changes of the original behavior, so they have
no legacy counterpart.
\begin{itemize}
\item \textbf{The high-level API}: in every language, a few functions that
  read and write whole records and annotation files as arrays, in one call,
  with named options and errors reported in the language's way
  (\cref{part:languages}).
\item \textbf{An output directory} (\code{wfdbx\_set\_output\_dir},
  \code{WFDBX\_OUTPUT\_DIR}, and the \code{dir} option of the high-level
  writers), in which output files are created, and an input directory
  (\code{wfdbx\_set\_input\_dir}, \code{WFDBX\_INPUT\_DIR}); see
  \cref{chap:files}.
\item \textbf{A record checker}, \code{check\_record(record, annotators)},
  which checks a record's header, signal files and annotation files, and
  reports each problem it finds with a severity (error, warning or
  information): missing or inconsistent signal specifications, ADC
  resolutions and ranges that the storage format cannot hold, samples outside
  the ADC range, checksum and length mismatches, problematic names,
  inconsistent segments of multi-segment records, and annotations that are
  out of order, undefined or after the end of the signals. Its report,
  including the text of every message, is identical in all seven
  implementations. It looks for problems similar to those found by Benjamin
  Moody's separate C program \cmd{wfdbcheck}, but was written independently.
\item \textbf{Fast reading and writing.} The high-level functions read and
  write whole blocks of samples, and \code{write\_annotations} encodes a whole
  MIT annotation file at once. The multi-frame readers added to the library
  functions (\code{getframes} in Python, R and MATLAB/Octave, \code{getvecs}
  in C++, Rust, Julia and Kotlin) decode whole blocks as well, including
  multi-frequency and multi-segment records. The samples and files are
  exactly those of the frame-by-frame functions.
\item \textbf{Signal names without surrounding spaces.} Some headers end a
  signal's description with spaces (\code{"ABP "}). The high-level readers
  give the names without white space at either end, and \code{read\_record}
  finds a channel by that name; the library functions and the tools keep the
  descriptions as the C library does.
\item \textbf{The header and the start time as numbers.} \code{read\_header}
  returns the header of a record without opening its signal files: the
  sampling and counter frequencies, the base counter value, the number of
  samples, the signal specifications, the info strings, the metadata, and the
  start time and date. \code{read\_header} and \code{read\_record} give the
  start time as text (\code{base\_time}), in milliseconds after midnight
  (\code{base\_time\_ms}, which keeps the milliseconds of a start time such
  as \code{14:00:03.500}), and the start date as year, month and day
  (\code{base\_date}).
\item \textbf{Write errors and frame checks in the high-level writers.}
  \code{write\_record} and \code{write\_annotations} raise an error
  (\code{write error on record REC}, \code{write error on annotator ANN for
  record REC}) when a file cannot be written completely, or when
  \code{putann} reports an error (version 1.2.2). \code{write\_record}
  also refuses, before creating any file, frame sizes that are not integers
  of at least 1 (\code{samps\_per\_frame must hold integers >= 1}) and
  signals whose number of columns is not the number of samples in a frame
  (\code{signals has N columns, but samps\_per\_frame adds up to M}).
\item \textbf{The high-level writers check every argument before touching
  any file, and leave the existing files intact after a failure} (version
  1.2.2). \code{write\_record} refuses an array of signals that is not 1-D or
  2-D (\code{signals must be a 1-D or 2-D array}), baselines, ADC zeros and
  ADC resolutions that are not whole numbers within the range of a 32-bit
  integer (\code{baselines must be whole numbers within the range of a
  32-bit integer}), and a record name, signal names, units or
  comments that cannot be written as header text (in Python, a character
  that is not Latin-1). \code{write\_annotations} refuses an array of times
  (or of codes, subtypes, channels or numbers) that is not 1-D
  (\code{samples must be a 1-D array}), annotation codes outside 0 to 58
  (\code{annotation code 61 is reserved (59 to 63 are the control codes of
  the annotation file format)}, \code{annotation code 64 is out of range (0
  to 58)}: the codes 59 to 63 are the control codes \code{SKIP}, \code{NUM},
  \code{SUB}, \code{CHN} and \code{AUX} of the MIT format, and an annotation
  written with one of them corrupts the file), and the times $-2^{63}$ and
  $2^{63}-1$ (\code{annotation time T cannot be written (annotation times
  must be integers from -9223372036854775807 to 9223372036854775806)}),
  which \code{putann} refuses (\cref{sec:known-c}); every other time can be
  written. Both writers then write their files to temporary files in the
  directories of the destinations (\file{.NAME.XXXXXXXX.tmp}), which are
  renamed over the destinations once all of them are complete, and removed
  after a failure, so that a failure never leaves a truncated header, signal
  file or annotation file in place of an existing one. A destination that is
  not a regular file, such as a device, is written directly; for a symbolic
  link, the file it leads to is replaced, and the link is kept. The files
  written are unchanged, byte for byte, and keep the permissions of the files
  they replace. (In legacy mode, an annotation file that must be sorted is
  renamed over its destination before the \cmd{sortann} program sorts it
  there.) The library functions (\code{putann}, \code{putvec},
  \code{newheader}, \ldots) and the tools keep the behavior of C.
\item \textbf{Read errors in the high-level readers.} \code{read\_record}
  raises an error (\code{can't read signal data of record REC (getvec status
  S)}) when a signal file ends before the length its header gives, or cannot
  be read, instead of returning the samples read so far. A read failure is
  tracked separately from the checksum: when a signal file ends in the frame
  at which the checksum of a signal is verified, for which \code{getvec}
  returns $-4$ (checksum error) instead of $-3$, the high-level readers
  still raise the read error (\code{getvec status -3}) and never return an
  incomplete frame, and \code{check\_record} reports the file as shorter
  than its header says, without a checksum error (version 1.2.2).
  \code{read\_record} allocates at most $2^{26}$
  samples in advance, whatever length the header gives, and grows its
  result as the samples are read, so that a header that gives far more
  samples than its files hold ends in the read error, not in an attempt to
  allocate memory for them. \code{read\_annotations} checks
  \code{sampfrom} and \code{sampto} as \code{read\_record} does (whole
  numbers, \code{sampfrom} not negative), before the file is opened, and
  closes the file whatever happens. A checksum error alone (every sample
  read, but their sum differs from the header's checksum) is a warning
  instead (\cref{sec:errors}).
\item \textbf{Header metadata} (version 1.2.0): the functions
  \code{read\_metadata}, \code{get\_metadata}, \code{set\_metadata},
  \code{remove\_metadata} and \code{update\_metadata} read and change the
  \code{Key: value} info strings of a header without changing anything else
  in it, \code{read\_header} returns them, and \code{write\_record} writes
  them (\cref{chap:metadata}). The info strings are those that
  \code{getinfo} reads, after the record line found as the header parser
  finds it, in every mode (\cref{sec:record-line-divergence}). The library
  functions \code{putinfo} and \code{setinfo}, and the tools, are unchanged,
  and the C library reads the changed headers as before.
\item \textbf{The high-level functions leave the annotation time resolution
  and the getvec mode as they were} (version 1.2.0). Opening an EDF record
  sets the time resolution of output annotators (\code{setafreq}) to the
  record's sampling frequency when none is set, and adds
  \code{WFDB\_HIGHRES} to the getvec mode (\code{setgvmode}), as in C, and
  \code{wfdbquit} resets neither. \code{read\_record}, \code{read\_header},
  \code{read\_annotations}, \code{write\_record}, \code{write\_annotations}
  and \code{check\_record} restore both previous values (\code{getafreq},
  \code{getgvmode}) when they return, also after an error, so that
  \code{read\_header("100s.edf")} changes nothing for a later
  \code{write\_annotations} or \code{getvec}; the \code{fs} argument of
  \code{write\_annotations} applies to its own file only. The library
  functions keep the behavior of C (\cref{sec:known-c}).
\item \textbf{Library functions} for these settings: \code{wfdbx\_set\_legacy},
  \code{wfdbx\_get\_legacy} and \code{wfdbx\_legacy\_args} (legacy mode),
  \code{wfdbx\_time\_needs\_freq}, \code{wfdbx\_write\_errors} and
  \code{wfdbx\_note\_write\_error}, the functions of the input and output
  directories and of network access, and \code{wfdbx\_input\_path} and
  \code{wfdbx\_output\_path} (version 1.2.0), which give the path under
  which a program opens a file that it reads or creates itself
  (\cref{sec:tool-files}).
\end{itemize}
The high-level functions behave identically in all seven implementations,
in legacy mode too, since they are not part of the original package.

\section{Corrections of version 1.2.0}
\label{sec:corrections-1-2-0}

Version 1.2.0 changes the behavior of WFDB-x in three points, which were
errors of the earlier versions:
\begin{itemize}
\item \textbf{The files that a tool reads or creates itself follow the input
  and output directories} (\cref{sec:tool-files}). The \opt{-i} file of
  \cmd{wrsamp}, \cmd{edf2mit}, \cmd{wav2mit}, \cmd{a2m} and \cmd{ad2m}, the
  input files of \cmd{ahaecg2mit} and the \opt{-S} script of \cmd{xform} are
  read from the input directory; the WAV file of \cmd{mit2wav}, the EDF file
  of \cmd{mit2edf}, the master header of \cmd{snip -m} and the \file{.des},
  \file{.ecg}, \file{.key} and \file{.txt} files of \cmd{parsescp} are
  created in the output directory; \cmd{wfdbdesc} measures a signal file of
  unknown length in the header's directory. They were opened relative to the
  current directory: with \code{WFDBX\_OUTPUT\_DIR} set, \cmd{mit2wav} left a
  WAV header without samples in the current directory and the samples
  without a header in the output directory. In legacy mode, input files are
  opened as named, and the output directory is a prefix of every output
  name. New library functions: \code{wfdbx\_input\_path} and
  \code{wfdbx\_output\_path} (\cref{sec:library-additions}).
\item \textbf{\code{osigopen(record, NULL, 0)}} returns 0 in every language
  (\cref{sec:crashes}); the Python implementation raised an
  \code{IndexError}.
\item \textbf{The high-level functions leave the annotation time resolution
  and the getvec mode as they were} (\cref{sec:additions}). After
  \code{read\_header("100s.edf")}, or any high-level function that opens an
  EDF record, a later \code{write\_annotations} in the same process no longer
  records the EDF record's sampling frequency as the time resolution of its
  file, and the getvec mode no longer stays in high resolution. The library
  functions keep the behavior of C.
\end{itemize}

\section{Corrections of version 1.2.2}
\label{sec:corrections-1-2-2}

Version 1.2.2 corrects the following points, in all seven implementations.
The library functions and the tools keep the behavior of C, except for
\code{getinfo} and the URL schemes.
\begin{itemize}
\item \textbf{Truncated records.} A signal file that ends before the length
  that its header gives, or that cannot be read, is a read error of
  \code{read\_record} (\code{can't read signal data of record REC (getvec
  status -3)}) even when the checksum of a signal is verified in the
  incomplete frame. \code{getvec} returns the checksum status ($-4$) for that
  frame, which replaced the end of file ($-3$), so \code{read\_record}
  returned the incomplete frame as data: a format-16 file holding the
  samples 1 and 2, whose header gives 3 samples, was read as 1, 2, 257, with
  a checksum warning. The high-level readers now track a read failure
  separately from the checksum, and never return an incomplete frame;
  \code{check\_record} reports such a file as shorter than its header says.
  A checksum error alone is still a warning (\cref{sec:additions}).
\item \textbf{\code{read\_record} allocates no more than it can read.} Its
  result was allocated from the length that the header gives (a header
  giving $10^{12}$ samples for a file of 4 bytes attempted to allocate
  terabytes); at most $2^{26}$ samples are now allocated in advance.
\item \textbf{The high-level writers check everything before touching any
  file, and leave the existing files intact after a failure}
  (\cref{sec:additions}). \code{write\_record} refuses, before any file is
  opened, arrays that are not 1-D or 2-D, baselines, ADC zeros and ADC
  resolutions that are not whole numbers within the range of a 32-bit
  integer, and text that cannot be written in a header (in Python,
  \code{units="}$\mu$\code{V"} left the existing header empty).
  \code{write\_annotations} refuses arrays of times that are not 1-D (a
  2-D array truncated the existing file), the annotation codes 59 to 63 and
  any code outside 0 to 58, and the times $-2^{63}$ and $2^{63}-1$, which
  were dropped while success was reported. An error that \code{putann}
  reports is an error of \code{write\_annotations}. Both writers write
  through temporary files that replace the destinations only when all of
  them are complete, keeping their permissions and symbolic links.
\item \textbf{\code{read\_annotations} checks \code{sampfrom} and
  \code{sampto}} as \code{read\_record} does (a fraction, NaN or an
  infinite value is an error, before the file is opened: \code{sampfrom
  must be a whole number}, \code{sampfrom must not be negative}), and closes
  the annotation file whatever happens.
\item \textbf{The record line of a header} is found by \code{getinfo} as
  the header parser finds it, except in legacy mode, and by the metadata
  functions in every mode (\cref{sec:record-line-divergence}): for a header
  whose first line is empty or blank, the comments before the record line
  were returned as info strings (in \code{comments} and in the notes of
  \cmd{wfdbdesc}).
\item \textbf{A remote file that could not be read is requested again}
  (\cref{sec:network}). A download that failed (a server error, a timeout, a
  file not yet published) was remembered until the process ended, so the
  file could not be read even after the server recovered; only the files
  that have been read are now kept.
\item \textbf{Only \code{http://} and \code{https://} URLs}
  (\cref{sec:network}). A name that begins with another URL scheme
  (\code{ftp://}, \code{sftp://}, \code{file://}, \ldots) was opened as a
  local path (\code{ftp://h/db/100} as the directory \file{ftp:}, then
  \file{h/db/100}); it is now refused in normal mode, with or without
  network access, with the message \code{WFDB-x: can't open NAME: only
  http:// and https:// URLs can be read (mount the server as a local folder,
  or download the files)}, and so are an input directory, a file named in a
  header, the files that a tool reads or creates itself, and an output
  directory with such a scheme. Legacy mode is unchanged.
\item \textbf{\cmd{time2sec -r} needs the sampling frequency of the record}:
  a record whose header gives none, or cannot be read, is an error (exit
  status 2: \code{time2sec: sampling frequency of record RECORD is not
  specified}). It printed \code{-9223372036854775808} with the exit status 0
  (a division by the unspecified frequency, 0); the C tool converted the time
  at 250\,Hz, as legacy mode still does (\cref{sec:no-defaults}).
\item \textbf{\code{write\_annotations} requires a label for each
  annotation}, in every language: with one annotation or more and neither
  symbols nor codes, it is an argument error (\code{symbols or codes must be
  given}; in R, \code{symbol or code must be given}), before any file is
  touched. C++, Rust and Julia wrote such annotations with the code 0,
  which is not an annotation code; Python, R and Kotlin refused even an
  empty list, which now writes an empty file, as in the other languages.
\item \textbf{\code{write\_record} removes a trailing \file{.hea}} from the
  record name, as the readers do: \code{write\_record("x.hea", \ldots)}
  writes \file{x.hea} and \file{x.dat}. Python failed (\code{can't create
  signal file for record x.hea}), and the other languages named the signal
  file \file{x.hea.dat}.
\item \textbf{Metadata changes} keep the permissions of the header, whatever
  its file name, and their temporary file is never more open than the header
  (\cref{sec:metadata-writing}); corrections in R, MATLAB/Octave and
  Kotlin.
\item \textbf{Licenses.} The file \file{NOTICE} (the licenses of
  \wfdbx{}, and the BSD-3-Clause notice of libFLAC, whose stream decoder the
  FLAC decoders translate) is included in the package of every language,
  with the license files (\cref{sec:license}). The R package declares
  \code{License: file LICENSE}, and in MATLAB/Octave the dispatcher of the
  \cmd{wfdbx} command (\file{+wfdbx/+internal/cli\_main.m} and
  \file{app\_run.m}), which only the tools use, is under the GPL, as the
  launchers of the other languages are.
\item \textbf{C++ package.} The installation provides
  \file{wfdbxConfigVersion.cmake}, so that
  \code{find\_package(wfdbx 1.2.2 CONFIG REQUIRED)} checks the version, and
  the \code{wfdbx} target requires C++17 of the programs that link to it.
\end{itemize}

\section{Legacy mode}
\label{sec:legacy}

Legacy mode makes \wfdbx{} behave exactly as the original package, for
reproducing old results. The default is the behavior described in the rest
of this manual. It is turned on in three ways:
\begin{itemize}
\item \textbf{Tools}: \opt{--legacy} as the \emph{first} argument of any tool
  (\cmd{rdsamp --legacy -r 100s -t 10}; in Python \cmd{wfdbx rdsamp --legacy
  -r 100s -t 10}). It is removed before the options are parsed, and the tool
  then behaves exactly as the original one: same options (the options added
  by \wfdbx{}, such as \opt{-F} in \cmd{rdann} or \opt{-R} in \cmd{wrsamp},
  are unrecognized), outputs, messages, help texts and exit statuses.
  \opt{--legacy} anywhere else is not recognized.
\item \textbf{Library}: \code{wfdbx\_set\_legacy(1)} turns legacy mode on,
  and \code{wfdbx\_set\_legacy(0)} off; each returns the previous setting,
  and \code{wfdbx\_get\_legacy()} returns the current one. The setting applies
  to the whole library. \code{wfdbx\_legacy\_args} is the function that the
  tools call to recognize and remove \opt{--legacy}.
\item \textbf{High-level API}: the \code{legacy} option of \code{read\_record},
  \code{read\_header}, \code{read\_annotations}, \code{write\_record},
  \code{write\_annotations} and the related functions turns the library's
  legacy mode on for the duration of the call, and restores the previous
  setting afterwards. The functions also follow the original rules when the
  library is in legacy mode.
\end{itemize}
No environment variable turns legacy mode on.

In legacy mode, every library difference of this chapter behaves as in the
original: a header without a frequency means 250\,Hz; a missing ADC
resolution is inferred from the format or taken as 12 bits; the time
conversions use 1 sample per second when no frequency is known;
\code{setifreq} reports \code{no open input record}; the unit conversions
assume a gain of 200 (and a baseline of 0 for invalid signal numbers); files
are found through the WFDB search path, without network access
(\cref{sec:resolve}); out-of-order annotation files are sorted by running
\cmd{sortann}, which loses annotations as the original does; an output file
may be one of the input files; \code{getinfo} takes the first line that does
not begin with \code{\#} for the record line; the whole record name is
checked when a header is written, and headers name records and files as
given. The input directory and \code{WFDBX\_NETWORK} are not used, no remote
file is read (every name containing \code{://} is refused), and the tools
open the files that they read themselves as named; the output directory is
used, also for the files that the tools create themselves. The tools exit
with the original status after a write error, and print no new messages. In \code{write\_record}, legacy mode
gives a gain of 200 and makes the units optional; its default ADC
resolution is the storage width of the format (8 bits for formats 80 and
508, 10 for 310 and 311, 12 for 212, 24 for 24 and 524, 32 for 32) and 16
bits for every other format, which differs from \cmd{wrsamp --legacy} only
for format 8 (12 bits there, as in the original \cmd{wrsamp}).

Legacy mode is tested against the unmodified C package by the legacy suite
of the conformance tests (886 cases, each run with \opt{--legacy}), and the
library's legacy values by the unit tests.

\begin{note}
Legacy mode is kept for compatibility only. The values it assumes are
exactly the silent defaults that \wfdbx{} removed, so results produced with
it carry the same risks: a time, interval, resampled signal or physical value
computed from an assumed frequency, gain, resolution or unit is wrong if the
assumption is, and nothing in the output says so.
\end{note}

\section{Errors and warnings}
\label{sec:errors}

\paragraph{The library functions} report errors as the C library does, in
every language: a function returns an error status (a negative number, a
null pointer, \ldots, as its C counterpart does), and the library prints a
message on the standard error output, such as \code{init: can't open header
for record 100}. \code{wfdbquiet()} stops the messages from being printed,
and \code{wfdbverbose()} prints them again; \code{wfdberror()} returns the
last message (before any error, \code{WFDB library version 10.7.0 (Jun 10
2022).}), and \code{wfdb\_error(format, \ldots)} reports a message as the
library does. The messages are those of the C library, byte for byte, so
that the tools print what the C tools print.

\paragraph{The high-level functions} report errors in the usual way of each
language, with a message that says which operation failed, such as
\code{can't read record 100s}; the library's own message, if any, is printed
too, unless \code{wfdbquiet()} was called, and returned by
\code{wfdberror()}.
\begin{itemize}
\item Python raises \code{OSError} (\code{IOError}) when a file cannot be
  read or written, \code{ValueError} for an invalid argument,
  \code{TypeError} for an argument of the wrong type, and
  \code{FileNotFoundError} when an output directory (\code{dir}) does not
  exist.
\item C++ throws \code{wfdbx::Error}, or, in the metadata functions (and
  in \code{write\_record}, for invalid metadata), one of its subclasses
  \code{ArgumentError} and \code{FileError} (in the namespace
  \code{wfdbx}).
\item Rust returns a \code{Result} whose error type is
  \code{wfdbx::record::Error}.
\item Julia throws a \code{WfdbError}, or, for some invalid arguments of
  the metadata functions and of \code{check\_record}, an
  \code{ArgumentError}.
\item R signals an R error (in the metadata functions, of class
  \code{wfdbx\_argument\_error} or \code{wfdbx\_file\_error}).
\item MATLAB/Octave throws an error, with an identifier such as
  \code{wfdbx:read\_record}.
\item Kotlin throws a \code{WfdbException}, or, for an invalid argument of
  the metadata functions, the writers and \code{readAnnotations}, its
  subclass \code{WfdbArgumentException}.
\end{itemize}
The chapters of \cref{part:languages} list the errors of each function.

\paragraph{The checksum warning.} When a high-level reader has read every
sample of a signal but their sum differs from the checksum of the header,
it returns every frame, the last one included, with the warning
\begin{lstlisting}[style=text]
checksum error in record REC: the signal data do not match the checksums in its header
\end{lstlisting}
in each language's usual way: Python \code{warnings} (a
\code{RuntimeWarning}), R \code{warning()}, Julia \code{@warn},
MATLAB/Octave \code{warning()} (identifier \code{wfdbx:checksum}), and the
standard error output in C++, Rust and Kotlin (also silenced by
\code{wfdbquiet()}; in Kotlin, an \code{onWarning} callback can receive it
instead). The \code{checksum\_warning} option (\code{checksumWarning} in
Kotlin) turns it off, and the record's \code{checksum\_ok} field is false in
either case, so that a program can test for it. The library function
\code{getvec} reports such an error as in C: \code{getvec: checksum error in
signal N}, and the status $-4$ for the last frame; the tools, such as
\cmd{rdsamp}, therefore stop before printing that frame, as the C tools do.
A signal file that ends early is a read error of the high-level readers,
not a checksum warning, even when the checksum is verified in its last,
incomplete frame (\cref{sec:additions}).

\section{Known C behaviors that WFDB-x reproduces}
\label{sec:known-c}

\wfdbx{} reproduces the C package exactly, including behaviors that may look
surprising; where the documentation of the C package and its code disagree,
\wfdbx{} follows the code. The file \file{docs/KNOWN\_C\_BEHAVIORS.md} lists
them, and says for each whether the official documentation describes it.
Among them:
\begin{itemize}
\item \textbf{Checksums.} \code{getvec} returns $-4$ for the frame at which a
  checksum mismatch is found, so that programs that loop while
  \code{getvec(\ldots) >= 0}, such as \cmd{rdsamp}, stop before printing the
  last frame.
\item \textbf{\code{getinfo} and the header parser disagree on where the
  record line is.} For a header whose first line is empty or blank, the C
  function \code{getinfo} takes that line for the record line, and returns
  the comments between it and the real record line as info strings. This
  is reproduced in legacy mode only: \wfdbx{}'s \code{getinfo} finds the
  record line as the header parser does (\cref{sec:record-line-divergence}).
\item \textbf{\code{putann} writes reserved annotation codes, and drops
  the extreme times.} \code{putann} does not check the annotation code: an
  annotation written with one of the codes 59 to 63, the control codes of
  the MIT format, is read back as a control word, which corrupts the rest of
  the file, and a code above 63 loses its high bits. An annotation at the
  time \code{WFDB\_TIME\_MIN} or \code{WFDB\_TIME\_MAX} ($-2^{63}$ or
  $2^{63}-1$) is refused with \code{putann: time overflow in annotation file
  N} and the status $-1$, and nothing is written for it. The library
  function \code{putann} of \wfdbx{} does the same; the high-level
  \code{write\_annotations} refuses such codes and times before writing
  anything (\cref{sec:additions}).
\item \textbf{Modification labels.} \code{annopen} consumes the \code{NOTE}
  annotations at time 0 (subtype 0) at the beginning of an annotation file,
  so \code{getann} does not return them; a file whose table redefines
  mnemonics makes every annotation file written later in the same process
  repeat the table.
\item \textbf{The time resolution note.} A program that calls
  \code{setafreq} with a positive value (as \cmd{wrann} does) writes a
  \code{\#\# time resolution:} note at the beginning of its output files.
\item \textbf{Duplicates at time 0.} An annotation at time 0 that follows
  another at time 0 counts as in order, so that annotations written twice at
  a time-0 location are not sorted, and both are kept.
\item \textbf{State that survives \code{wfdbquit}}: the counter frequency
  and base counter value, the output annotation time resolution, the buffer
  sizes, and the mnemonic and description tables are not reset.
\item \textbf{Opening an EDF record} sets the time resolution of output
  annotators to the record's sampling frequency, when \code{setafreq} has not
  set one, and adds \code{WFDB\_HIGHRES} to the getvec mode; \code{wfdbquit}
  resets neither, so that an annotation file written later by the same
  process begins with a time-resolution note (\code{\#\# time resolution:
  360} after \code{100s.edf}), and later records are read in high
  resolution. The library functions do the same; the high-level functions
  restore both settings (\cref{sec:additions}).
\item \textbf{Tools.} \cmd{sortann} keeps only the last of several
  annotations with the same time, \code{num} and \code{chan}; \cmd{wrann}
  turns an unrecognized mnemonic into a \code{NOTE} whose aux string starts
  with the mnemonic, and keeps the subtype, \code{chan} and \code{num} of the
  previous line when a line omits them; \cmd{rr2ann} repeats the previous
  interval for an empty line; \cmd{revise} rejects an old-style header of
  exactly 32 signals; \cmd{edf2mit} adds stray bytes to a 16-character signal
  label; \cmd{ahaecg2mit} writes the annotations of \file{.ecg} input one
  sample early; \cmd{mit2edf} uses the current date and time for a record
  that has no start time.
\item \textbf{Documentation and code disagree} on a few points, and \wfdbx{}
  follows the code: \code{wfdbquit} does not reset the counter frequency;
  the variable that turns sorting off is \code{WFDBANNSORT}, not
  \code{WFDBNOSORT}; \cmd{wrsamp}'s \opt{-c} has no effect and \opt{-l} is
  obsolete; its default field separators are space, tab, CR, LF and comma;
  \cmd{edf2mit} and \cmd{wav2mit} name their default output records after
  the input file name.
\item \textbf{Build-specific values.} \code{wfdbversion}, \code{wfdbldflags},
  \code{wfdbcflags} and \code{wfdbdefwfdb} return the values of a release
  build of 10.7.0 installed in \file{/usr/local}: \code{"10.7.0"},
  \code{""}, \code{""}, and the default WFDB path of the C package,
  \code{". /usr/local/\allowbreak database
  http://\allowbreak physionet.org/\allowbreak physiobank/\allowbreak database"};
  \code{wfdbdefwfdbcal} returns \code{"wfdbcal"}.
\end{itemize}

\paragraph{Behaviors that depend on the platform.} The reference results were
produced on x86-64 Linux with glibc, where \code{char} is signed, \code{long}
is 64 bits, and \code{WFDB\_Time} is 64 bits. \wfdbx{} reproduces that
platform's behavior where C's would differ elsewhere; for example,
floating-point operations are never contracted into fused multiply-adds
(which some compilers do by default, and which can change results such as
the baselines that \cmd{edf2mit} computes). One case remains: when an output
file is also an input file in legacy mode (\cref{sec:output-input}), what is
read depends on how far the runtime reads ahead, as in C; the R
implementation on macOS reads further ahead than glibc, and fails one legacy
conformance case for that reason.
